\section{Conclusion}
This report has addressed the growing importance of Incident Response (IR) in the context of today’s increasingly complex and evolving threat landscape. As cyberattacks become more frequent, targeted, and sophisticated, the necessity for organizations to implement well-defined incident response capabilities is a must to guarantee the continuity of operations and the minimization of damage.

Through the examination of leading IR frameworks, we have assessed the operational strengths and limitations of methodologies such as NIST, SANS, and others. This report has specially focused on the NIST framework, from which all its phases have been analized in depth. 

Moreover, the report has analyzed and emphasized in the close relationship that Threat Intelligence and Incident Response must have, specially for the analysis of cyber threat actor motivations, behaviours, TTPs an the integration of threat intelligence within incident response planning.

Finally, as a result of all the work, the report provides a set of 3 example playbooks to existent threats. By studying real-world use cases, some response plans have been develop and they can actually be used as a base to develop playbooks to real existent threats. These examples, have mapped the NIST IR framework into theoretical practise. 

In summary, this project has underlined the critical role of incident response in modern cybersecurity. By combining threat intelligence, structured frameworks, real-world observations, and targeted recommendations, organizations can significantly strengthen their ability to detect, respond to, and learn from security incidents in a continuously learning and hardening loop.

